\section{Experiments}
\label{sec:experiments}

Because \method's current release operates at inference time only
(Section~\ref{sec:method:training}), our experiments focus on
validating the \emph{primitives} and the \emph{evaluation
methodology}. We establish three results:

\begin{itemize}
\item[\textbf{E1}] The language-memory schema round-trips losslessly
      through text rendering and parsing
      (Section~\ref{sec:exp:roundtrip}).
\item[\textbf{E2}] The persistent cross-session store retrieves the
      correct episode in a synthetic benchmark with non-trivial
      precision even using only a toy embedder
      (Section~\ref{sec:exp:retrieval}).
\item[\textbf{E3}] The Mem-0 integration successfully inserts
      retrieved facts into the planner prompt without breaking the
      original prompt structure
      (Section~\ref{sec:exp:integration}).
\end{itemize}

We additionally present the design of \textbf{MNEMO-Bench}
(Section~\ref{sec:exp:bench}), a unified harness consolidating four
existing memory benchmarks, along with the projected evaluation
protocol for full VLA training experiments
(Section~\ref{sec:exp:plan}).

\subsection{E1: Language Memory Round-Trip}
\label{sec:exp:roundtrip}

We construct memories with 3 subtasks, 1 cross-session fact, and
1 failure event, render them via \texttt{Memory.to\_text()}, and
parse them back via \texttt{parse\_memory\_text()}. Every field
(task, session id, subtask status, details, source session,
confidence, failure intervention) is recovered exactly. Across 4
unit tests covering the schema and parser modules, the round-trip
success rate is 100\%.

\subsection{E2: Synthetic Retrieval Benchmark}
\label{sec:exp:retrieval}

\paragraph{Setup.}
We seed a persistent store with 20 synthetic episodes across 5 task
families (kitchen, workshop, laundry, garden, office), each family
containing 4 episodes with distinct narratives. Queries are
constructed by paraphrasing target episodes. We measure Recall@1
(is the correct episode the top hit?) and Recall@3.

\paragraph{Embedders.}
We compare (a) \texttt{HashingEmbedder(dim=64)}, a zero-dependency
hasher, and (b) \texttt{SentenceTransformerEmbedder} using the
\texttt{all-MiniLM-L6-v2} checkpoint (when available in the
environment).

\paragraph{Results.}
Table~\ref{tab:retrieval} reports Recall@$K$ on 20 paraphrased
queries, all executed on a single CPU in under one second.

\begin{table}[h]
\centering
\small
\begin{tabular}{lccc}
\toprule
Embedder & R@1 & R@3 & R@5 \\
\midrule
\texttt{HashingEmbedder}(dim=128) & 75.0\% & 90.0\% & 95.0\% \\
\texttt{SentenceTransformer}(MiniLM) & \multicolumn{3}{c}{\textit{optional extra}} \\
\bottomrule
\end{tabular}
\caption{Synthetic cross-session retrieval Recall@$K$. The hashing
embedder already surpasses random (1/20 = 5\%) by a large margin;
the sentence-transformer variant is expected to reach near-ceiling
performance. Results are reproducible via
\texttt{examples/retrieval\_benchmark.py}.}
\label{tab:retrieval}
\end{table}

The hashing embedder achieves Recall@1 of 75\% and Recall@5 of
95\%---a 15$\times$ improvement over the 5\% random baseline on a
20-way classification. This surprisingly strong result confirms that
the retrieval infrastructure is functioning correctly and that even
a dependency-free deployment can offer useful cross-session recall
for deployments where bringing in a transformer is impractical.

\subsection{E3: Mem-0 Prompt Integration}
\label{sec:exp:integration}

We simulate a Mem-0 planner prompt (two-message list with system and
user roles, user content containing \texttt{<global\_task>} and
\texttt{<finished\_subtasks>} text entries) and verify that
\texttt{inject\_into\_prompt} (i) inserts the retrieved
\texttt{<cross-session memory>} block immediately before the
\texttt{<finished\_subtasks>} marker, (ii) preserves all original
content, (iii) does nothing if the store is empty. All three
properties are validated by unit tests in
\texttt{tests/test\_integrations.py}.

\subsection{MNEMO-Bench: A Unified Memory Harness (Design)}
\label{sec:exp:bench}

A direct comparison of memory-aware VLAs has been blocked by the
four-way fragmentation of existing benchmarks. We propose
\textbf{MNEMO-Bench}, a unified evaluation harness that:

\begin{enumerate}
\item Ports the nine tasks of RMBench \cite{rmbench} unchanged.
\item Re-implements the three tasks of MemoryBench \cite{sam2act}
      inside the RoboTwin 2.0 \cite{robotwin2} simulator for unified
      infrastructure.
\item Selects 10 representative tasks from MIKASA-Robo
      \cite{mikasa} covering the 12 memory-type categories.
\item Selects 10 representative tasks from RoboCerebra
      \cite{robocerebra} for long-horizon evaluation.
\end{enumerate}

The harness reports a per-memory-type \textbf{radar chart}:
\emph{spatial}, \emph{sequential}, \emph{episodic},
\emph{temporal}, \emph{visual}, \emph{capacity}, and
\emph{cross-session}. A single scalar ``overall memory score'' is
defined as the geometric mean across memory types, encouraging
methods that are uniformly capable rather than specialists.

The proposed protocol fixes: 30 rollouts per task for statistical
significance; 50 training demonstrations per task; and a fixed
domain-randomization strength drawn from RoboTwin 2.0. Results are
reported under both \textit{Easy} (clean) and \textit{Hard} (fully
randomized) settings.

\subsection{Experiment Plan for Full VLA Training}
\label{sec:exp:plan}

Our follow-up work will evaluate the training objective of
Equation~\ref{eq:joint_loss} along three axes:

\begin{itemize}
\item \textbf{Ablation on persistent memory.} Trains the Mem-0
      executor unchanged; compares inference with vs.\ without the
      persistent store on RMBench $M(n)$ and on a new
      \emph{cross-session kitchen} task.
\item \textbf{Consistency loss impact.} Measures how much the
      consistency term of Equation~\ref{eq:joint_loss} reduces
      memory-action divergence (Mem-0 currently reports a
      manually-tuned $\lambda_{\text{classifier}} = 0.2$).
\item \textbf{Cross-embodiment transfer.} Trains on Aloha-AgileX and
      tests transfer to Piper / Franka / UR5 / ARX-X5 using small
      per-embodiment adapters.
\end{itemize}

Estimated compute budget: 8$\times$H100 for 5 weeks of pretrain +
fine-tuning, plus 2 Aloha-AgileX platforms for 1,000 real-world
demonstrations and 300 rollout evaluations.
